\documentclass[10pt,twoside]{article}
\usepackage[utf8]{inputenc}
\usepackage{amsmath}
\usepackage{amsfonts}
\usepackage{amssymb}
\usepackage[spanish,es-noshorthands]{babel}
\usepackage[T1]{fontenc}
\usepackage{lmodern}
\usepackage{graphicx,hyperref}
\usepackage{tikz,pgf}
\usepackage{multicol}
\usepackage{subfig}
\usepackage[papersize={6.5in,8.5in},width=5.5in,height=7in]{geometry}
\usepackage{fancyhdr}
\pagestyle{fancy}
\fancyhead[LE]{\includegraphics[height=12pt]{Images/logo-colegio.png} Cálculo $11^{\circ}$}
\fancyhead[RE]{}
\fancyhead[RO]{\textit{Germ\'an Avenda\~no Ram\'irez, Lic. U.D., M.Sc. U.N.}}
\fancyhead[LO]{}

\author{Germ\'an Avenda\~no Ram\'irez, Lic. U.D., M.Sc. U.N.}
\title{\begin{minipage}{.2\textwidth}
\includegraphics[height=1.75cm]{Images/logo-colegio.png}\end{minipage}
\begin{minipage}{.55\textwidth}
\begin{center}
Taller, Límites de funciones $\mathbb{R}$\\
Cálculo $11^{\circ}$
\end{center}
\end{minipage}\hfill
\begin{minipage}{.2\textwidth}
\includegraphics[height=1.75cm]{Images/logo-sed.png} 
\end{minipage}}
\date{}
\begin{document}
\maketitle
Nombre: \hrulefill Curso: \underline{\hspace*{44pt}} Fecha: \underline{\hspace*{2.5cm}}
\section*{Introducci\'{o}n}
\subparagraph*{Materiales:}
Regla, escuadra, calculadora, esferos o lápices de diferentes colores.
\begin{enumerate}
\item Grafica cada una de las siguientes funciones definidas en el conjunto de los números Reales:
\begin{enumerate}
\begin{multicols}{3}
\item $y=f(x)=2x+1$
\item $y=g(x)=x^{2}-4$
\item $y=h(x)=x^{3}-2x$
\end{multicols}
\end{enumerate}
\item En la siguiente recta numérica, escoge un par de unidades consecutivas y cada una divídelas en 10 partes
iguales. Coloca el número correspondiente a cada división. ¿Cuáles serían los números si cada unidad es dividida en 100 partes iguales?
\begin{center}
\begin{tikzpicture}[>=stealth,scale=1.2]
\draw[<->] (-2.5,0) -- (3.5,0);
\draw [thick] (-2,-.1) node[below]{$-2$} -- (-2,0.1);
\draw [thick] (-1,-.1) node[below]{$-1$} -- (-1,0.1);
\draw [thick] (0,-.1) node[below]{0} -- (0,0.1);
\draw [thick] (1,-.1) node[below]{1} -- (1,0.1);
\draw [thick] (2,-.1) node[below]{2} -- (2,0.1);
\draw [thick] (3,-.1) node[below]{3} -- (3,0.1);
\end{tikzpicture}
\end{center}
\item A continuación encontrarás dibujadas dos rectas. Traza perpendiculares según se te pide:
\begin{center}
\begin{tikzpicture}[>=stealth]
\draw[<->] (-2.5,0) -- (3.5,0);
\draw (-2,0)circle (2pt);
\end{tikzpicture}
\end{center}
\end{enumerate}
\end{document}
